%==============================================================
\section{Tracking a drifting baseline (constant-gain estimators)}\label{sec:tracking}
%==============================================================

This section summarises the drift-model study (``Problem~14''). Unlike Sections~\ref{sec:curvature}--\ref{sec:tv} there is no theorem about the
estimator here: the three deterministic formulas below are proved for a \emph{linear filter}, and their relevance to the point-process
estimators is \Nt{}.

\subsection{The model class and the target form}

\begin{definition}[piecewise-Lipschitz class $\mathcal P(\Delta,J,D;c_\mu,c_\eps)$]\label{def:drift}
The true path $\theta^*_t$, $t\in[0,T]$, satisfies (box) $\mu_t\in[\mu,c_\mu\mu]$, $\eps_t=1-n_t\in[\eps,c_\eps\eps]$; (variation) $\theta^*_t$ is Lipschitz with $\norm{\dot\theta^*_t}\le\Delta$ between at most $J$ jump times, each jump of size at most $D$. Path length $V_T\le\Delta T+JD$. Drift is measured in a fixed reference Fisher norm $\norm\cdot_{I(\theta_{\mathrm{ref}})}$.
\end{definition}
Rationale: smooth-only drift cannot represent regime changes and misleads step-size tuning (jump bias $\propto1/\eta$, smooth bias $\propto1/\eta^2$); piecewise-constant only excludes smooth seasonality and needs every epoch longer than the relaxation time $1/(\eta r)$; a bare bounded-variation budget cannot separate the exponents.
The target is the average tracking MSE of a constant-gain estimator with gain $\eta$ (per event):
\[
\underbrace{\eta\,\sigma^2(\eps)}_{\text{noise}}\ +\ \underbrace{\Delta^2/\eta^2}_{\text{lag}}\ +\ \underbrace{J D^2/(\eta T)}_{\text{jumps (event units, up to a factor }1/2)} .
\]

\begin{proposition}[metric equivalence along the path]\label{prop:metric}\Nt{}
Let $\kappa_I=\max_{\theta,\theta'\in\text{box}}\lambda_{\max}(I(\theta')^{-1/2}I(\theta)I(\theta')^{-1/2})$. For $K=1$ (exact $I$) with $\mu$ and $\eps$ each varying by $\times1.5$: $\kappa_I=19.5,22.5,25.6$ at $\eps_0=0.2,0.1,0.05$, essentially independent of $\eps_0$, whereas $\mathrm{cond}\,I\sim\eps^{-2}$ ($48,152,522$). For $K=2$ ($\beta=(1,2)$, Monte Carlo) $\kappa_I\approx27$--$30$. Converting drift to a fixed reference norm costs $\kappa_I^{1/2}\approx2$--$5$, not $\eps^{-2}$, for a $\times1.5$ box; it grows fast with the box ($\times2$ in both: $150$--$220$).
\end{proposition}

\subsection{The three terms for a linear filter}

Let the error obey $e_{t+1}=(1-\eta)e_t+\eta\xi_t-d_t$ with $\xi_t$ martingale-difference noise and $d_t=\theta^*_{t+1}-\theta^*_t$ (Fisher-preconditioned coordinates).

\begin{proposition}[noise]\label{prop:noise}\Pt{}
If $\Var\xi_t=\sigma_\xi^2$ then $\Var e=\eta\sigma_\xi^2/(2-\eta)\approx\eta\sigma_\xi^2/2$.
\end{proposition}

\begin{proposition}[sinusoidal drift]\label{prop:sin}\Pt{}
For the continuous filter $\dot y=\kappa(u-y)$ and $u=u_0+A\sin\omega t$, the steady error has amplitude $A\omega/\sqrt{\kappa^2+\omega^2}$ and mean square $\frac{A^2}2\frac{\omega^2}{\kappa^2+\omega^2}$; it is $\approx\Delta^2/\eta^2$ ($\Delta=A\omega$) when $\kappa\gg\omega$.
\end{proposition}

\begin{proposition}[alternating jumps]\label{prop:tanh}\Pt{}
Let $u$ alternate between two levels differing by $D$, switching every $L$ time units, and $\dot y=\kappa(u-y)$. The time-averaged squared error in steady state is
\[
\overline{e^2}=\frac{D^2\tanh(\kappa L/2)}{2\kappa L}\ \xrightarrow{\ \kappa L\gg1\ }\ \frac{D^2}{2\kappa L}=\frac{\rho D^2}{2\kappa},\ \ \rho=1/L;\qquad\xrightarrow{\ \kappa L\ll1\ }\frac{D^2}4 .
\]
\end{proposition}
\begin{proof}
At a switch the error jumps by $D$; periodicity and symmetry give the post-switch magnitude $e_0=D-e_0e^{-\kappa L}$, i.e.\ $e_0=D/(1+e^{-\kappa L})$. On a period $e(t)=e_0e^{-\kappa t}$, so $\frac1L\int_0^Le^2=e_0^2(1-e^{-2\kappa L})/(2\kappa L)=D^2\frac{1-e^{-\kappa L}}{1+e^{-\kappa L}}\frac1{2\kappa L}$.
\end{proof}
(Independently re-derived in this review; it agrees with the formula in the source note.)

\begin{proposition}[path-length bound]\label{prop:pathlen}\Pt{}
$e_t=-\sum_{s\ge0}(1-\eta)^sd_{t-s}$, so by Young's inequality $\frac1T\sum e_t^2\le\eta^{-2}\frac1T\sum d_t^2\le\Delta_{\max}V_T/(\eta^2T)$. Also $\sum e_t^2\le\norm e_\infty\norm e_1\le V_T\cdot V_T/\eta$, i.e.\ $\frac1T\sum e_t^2\le V_T^2/(\eta T)$.
\end{proposition}
\begin{remark}
The source note states the second bound with an extra factor $\frac12$ ($V_T^2/(2\eta T)$). The argument above only gives $V_T^2/(\eta T)$; the sharper constant is not derived in the note. Flagged in Section~\ref{sec:review}.
\end{remark}

\paragraph{Time-varying gain.} A per-event gain $\eta$ is a per-unit-time gain $\kappa_t=\eta\,r_t$ with $r_t=\mu_t/\eps_t$ the local event rate. The correct deterministic baseline is $\dot y=\eta r_t(\mu_t-y)$, not the constant-gain filter; it reproduces triangle-wave and ramp results to within $5\%$ (Table~\ref{tab:track}). The ``lag'' term $\Delta^2/\eta^2$ is the response of this filter and equals $\Delta^2/\eta^2$ only for slow drift and small relative amplitude.

\subsection{Evidence (all \Nt{}; causal sampling, $\hat\mu$ read at the last event at or before each grid time)}

Two constant-gain estimators were used: ``naive'' ($\hat\mu\mathrel{+}=\eta(r/m)(1/\hat\lambda-w)$) and ``forgetting Newton'' ($F\leftarrow(1-\eta)F+\Delta F_i$, $\hat\theta\mathrel{-}=F^{-1}g_i$, the constant-gain version of Section~\ref{sec:regret}); Problem~5's estimator verbatim ($\eta=0$) does \emph{not} track (the gain decays like $1/n$; bias$^2\approx0.19$--$0.20$ in four successive windows) and needs forgetting.

\begin{table}[h]
\centering\small
\begin{tabular}{@{}p{4.2cm}p{9.6cm}@{}}\toprule
quantity & finding\\\midrule
noise coefficient, $K=1$, $\alpha$ known & $\mathrm{MSE}\approx a\eta$, $a=\mu/(2\eps m)$ with exact $m$; measured/predicted $=1.09$--$1.18$ across $k=0.6$ to $8$; $a\eps^2/\mu^2=0.27$ ($k=0.6$) to $0.52$ ($k\ge4$); $k=2$: $0.44$\\
noise, $\alpha$ unknown & $a=r/(2\,\eff)$, $\eff=(m-\eps/\mu)/(1-2\eps+\eps\mu m)$ (checked against the Fisher formulas of Section~\ref{sec:fisherK1}); $\eps=0.2$: $22.2$ predicted vs $17.6$--$21.8$ measured; $K=2$: $23.2\eta$ vs $23.7\eta$; $K=3$ $(1,3,9)$: $18.7\eta$ vs $19.5\eta$\\
exact rescaling & $t\to t/c$, $\mu\to c\mu$, $\beta\to c\beta$ leaves the per-event recursion unchanged; only $k=2\mu/\beta$ matters\\
linear window & the filter formula holds while $a\eta/\mu^2\lesssim\theta_k$, $\theta_k\approx0.013$ ($k=2$), $0.015$ ($k=4$), $0.03$ ($k=0.6$); i.e.\ $\omega\lesssim\eta\lesssim\eta_{\max}=\theta_k\mu^2/a=2\theta_k\eps m\mu$; non-empty iff the drift is slow enough\\
transfer function ($K=1$, $\eps=0.4$, $A=0.3$, $P=1200$) & $\eta=0.004$: $H=0.768-0.428i$ ($\alpha$ known), $0.799-0.411i$ (fixed preconditioner), $0.814-0.518i$ (adaptive) vs filter $0.785-0.411i$; within $1$--$2$ s.e.\ except adaptive at $\eta=0.004$ (modulus $+9\%$)\\
decomposition of the error at $\eta=0.004$ & fundamental error$^2$/filter: fixed preconditioner $1.07,0.98,0.92$; adaptive $1.41,1.48,1.61$ for $A=0.2,0.35,0.5$; adaptive/fixed $=1.32,1.48,1.75$: an \emph{amplitude-dependent} excess of the data-driven preconditioner (absent for a fixed one, $\to1$ as $A\to0$); not derived analytically\\
$K\ge2$ ($\eps=0.4$, causal) & within the window the noise coefficient holds to $3$--$13\%$ and $H$ follows the filter to $2\%$ ($K=3$, $(1,3,9)$), $6\%$ ($K=2$, $(1,2)$), $21\%$ ($K=3$, $(1,2,4)$); not monotone in $\mathrm{cond}\,I$ ($755$ vs $1790$); outside the window ($a\eta/\mu^2\approx0.04$) modulus overshoots and noise rises $30$--$40\%$\\
other shapes ($K=1$) & ramp: lag $=(\text{slope}/r_{\mathrm{local}})/\eta$ within $15\%$; triangle wave: $0.0154,0.0066$ measured vs corner-aware filter $0.0161,0.0068$ (steady-state formula off by $2.3\times$); random walk: measured MSE $20$--$40\%$ above $a\eta+(q/r)/(2\eta)$\\
jump class & ratio measured/(time-varying filter): $1.00$--$1.10$ ($\alpha$ known), $1.02$--$1.43$ ($\alpha$ unknown) inside the window; $2.6$--$9$ where the window rule says $\eta$ is too large\\\bottomrule
\end{tabular}
\caption{Summary of the tracking experiments (source: the drift-model note, sections F1--F4 and R0--R7).}\label{tab:track}
\end{table}

\paragraph{Corrections already made in the source note.} (i) The first versions of the mechanism and $K=2$ scripts read $\hat\mu$ at the first event \emph{after} each grid time, which peeks ahead; the $K=2$ bias ratios $1.64,1.87$ are superseded. (ii) The claim $c_{\mathrm{lag}}\approx1.8$ is retracted (the mean-square ratio varied between $1.2$ and $1.9$ across seed sets for the same settings). (iii) Hall--Willett (arXiv 1409.0031, full text) assume no drift model: their regret depends on $\sum\norm{\lambda_{t+1}-\Phi_t(\lambda_t)}^2$ with a $1/\sqrt t$ step in the rate vectors, a different object from a parameter MSE.

\paragraph{Open.} $K\ge4$; $\eps<0.4$ for $K\ge2$ (needs $\eta\lesssim10^{-3}$, runs of $10^5$ time units); a closed form for the amplitude effect in the matrix ($\alpha$ unknown) case (the scalar case is derived in Section~\ref{sec:resolutions}); the window hypothesis (W) of Theorem~\ref{thm:track}; the full text of Piggott--Solo (ICASSP 2018, closed access) to check (1) whether a tracking bound is stated, (2) its parameter-variation model, (3) the forgetting factor, (4) the two-time-scale non-negativity mechanism versus projection, (5) the Laguerre basis versus the exponential dictionary.

%==============================================================
\section{Review of the whole programme}\label{sec:review}
%==============================================================

\subsection{Concordance of labels}

\begin{center}\small
\begin{tabular}{@{}p{5.6cm}p{7.5cm}@{}}\toprule
source-note label & this document\\\midrule
Problem~1 Lemmas 1,2,3 & Lemmas~\ref{lem:rankone},\ref{lem:expcurv},\ref{lem:lammin}\\
Problem~3 key identity, Lemma 1/1$'$, Lemma 4 & Prop.~\ref{prop:keyid}, Lemma~\ref{lem:moments}, Lemma~\ref{lem:expmoment}\\
Problem~3 Theorem 2, Prop.\ 3, Lemmas A,C & Thm.~\ref{thm:contraction}, Prop.~\ref{prop:gamma}, Lemma~\ref{lem:ellcontr}\\
Problem~10 Theorem A (and repaired Step 2) & Thm.~\ref{thm:A}, Lemma~\ref{lem:gammafact}\\
Problem~11 Theorem B & Thm.~\ref{thm:B}\\
Problem~11 Prop.\ 3; Lemma G, H; Theorems E, F; Cor.\ D/E$'$ & Prop.~\ref{prop:ridge}; Lemmas~\ref{lem:G},\ref{lem:H}; Thms.~\ref{thm:E},\ref{thm:F}; Cor.~\ref{cor:Eprime}\\
Problem~11 ``Theorem C'' (Lyapunov, region $2\mu>n^2B$) & superseded by Thm.~\ref{thm:E}; not reproduced\\
Problem~2 ``Theorem T'' (truncation) & Thm.~\ref{thm:trunc}\\
Problem~12 & Prop.~\ref{prop:tiers}\\
Problem~4 & Props.~\ref{prop:projdiag},\ref{prop:projbox}\\
Problem~5 Prop.\ 1, ``Theorem B'' (box) & Props.~\ref{prop:reg},\ref{prop:ratio}\\
Problem~7 Lemmas 7a--c, Theorem 7 & Lemmas~\ref{lem:score},\ref{lem:drift},\ref{lem:logdet}; Thm.~\ref{thm:cond}\\
Fisher concentration Lemmas F1,F2, ``Theorem F'' & Lemmas~\ref{lem:F1},\ref{lem:F2}; Thm.~\ref{thm:Fconc}\\
Exponential tails E1--E4; E5; ``Theorem C-exp'' & Lemmas~\ref{lem:E1},\ref{lem:E2},\ref{lem:E4}, Thm.~\ref{thm:E3}, Lemmas~\ref{lem:E5a}--\ref{lem:E5d}, Prop.~\ref{prop:E5}; Thm.~\ref{thm:Cexp}\\
Bootstrap ``Theorem C'' (polynomial tails) & superseded by Thm.~\ref{thm:Cexp}\\
Expected regret ``Theorem R'' & Thm.~\ref{thm:R}\\
Stationary-TV ``Theorem T'', ``Lemma F'' & Thm.~\ref{thm:Tprime}, Lemma~\ref{lem:LemmaF}\\
General-$K$ ``Theorem V'', ``Lemma S'' & Thm.~\ref{thm:V}, Lemma~\ref{lem:S}\\\bottomrule
\end{tabular}
\end{center}

\subsection{Errata found and corrected during the programme}

\begin{enumerate}[nosep]
\item \textbf{Thinning bug.} Several Monte Carlo scripts discarded rejected proposals instead of accumulating waiting time; affected the first versions of the Problem~2 tables, the Problem~3 exponential-tail check, the Problem~4 $K\ge2$ conditioning table and the first transient check in Problem~10. All were fixed and the affected numbers regenerated; the proved statements were not affected.
\item \textbf{Problem~2 erratum.} ``$\lambda_{\min}(I)$ does not shrink with $\eps$ for $K\ge2$'' was a range artefact ($\eps\ge0.05$ measures the middle tier). For $\eps\to0$, $\lambda_{\min}\propto\eps$ for every $K$.
\item \textbf{Diagonal clip.} The diagonal-metric projection is not non-expansive in the $F$-norm and inflates regret $\approx10\times$ at tight boxes; for $K=1$ the exact $F$-projection (Prop.~\ref{prop:projbox}) must be used.
\item \textbf{Theorem~A, Step 2.} Two inadequate justifications (the $x\ge1.5$ bound needs convexity of $h$; the denominator bound needs the exact factor $r\le\frac13$) were repaired (Lemma~\ref{lem:gammafact}).
\item \textbf{Initial-state gap in Corollary~E$'$ (and an overstatement found in the final audit).} The KL lemma compares laws with the same initial state; the stationary experiment has $S_0\sim\pi_\theta$ vs $\pi_{\theta'}$. Closed for $K=1$ via $\TV(\pi_\theta,\pi_{\theta'})$ (Thm.~\ref{thm:Tprime}); bypassed for all $K$ via ghost immigrants (Thm.~\ref{thm:V}). In addition, the notes' claim that ``given the pre-sample state $S_0$ no estimator is consistent'' is \emph{not} implied by the KL lemma (the law of $S_0$ depends on $\theta$; with $S_0$ ancillary the information is not $T\,I$); Theorem~\ref{thm:B} and Corollary~\ref{cor:Eprime} are restated accordingly.
\item \textbf{Theorem~7 / Lemma~7c.} The assumption ``gap leverage $x_i\le\bar x$ almost surely'' cannot hold ($\tr F^{-1}\Delta F\propto E_i\sim\mathrm{Exp}(1)$); replaced by the truncated Lemma~\ref{lem:logdet}.
\item \textbf{$|d(y)|$ constant.} The earlier bound $31.5/\alpha$ used $|N|\le2.5x^2$ without proof; replaced by $19/\alpha$ with a complete argument ($\mathrm{Re}\,N\in[-y^2/2,0]$); true supremum $1.1429/\alpha$.
\item \textbf{Drift study.} Peek-ahead sampling in two scripts and the retracted $c_{\mathrm{lag}}\approx1.8$ (Section~\ref{sec:tracking}).
\item \textbf{Constant in Lemma~\ref{lem:logdet}.} The first draft bounded $\tr(F^{-1}\Delta F)$ by $\bar\ell c^+E_i/\mu^*$; the correct bound is $\bar\ell E_i/(c^-\mu^*)$ (since $\lambda_\theta\ge c^-\lambda^*$). Harmless for the leading constant ($\bar x\to0$) but wrong as written; the K=2 pathwise audit could not detect it because the iterates never reached the extreme of the box.
\item \textbf{Test-script error (not a mathematical one).} A pointwise check of $|\Delta'|\le M_1\eps(k(e-1)+e)|G|$ initially failed because the code compared an $x$-derivative with a $t$-bound; corrected, the inequality holds.
\end{enumerate}

\subsection{Inconsistencies between the notes that remain}

\begin{enumerate}[nosep]
\item \textbf{Label clashes} (three ``Theorem B'', three ``Theorem C'', two ``Theorem T'', two ``Lemma F'', ``Lemma 3'' in two senses). Resolved here; the map for the notes is in \texttt{LABELS.md}.
\item \textbf{Drift note, stale paragraphs} \emph{(cleaned in the notes on this date; kept here as a record).} Section F2 still contains ``empirically the lag term should carry a constant $c_{\mathrm{lag}}\approx1.8\pm0.2$'' and the statement that the excess is ``independent of $\eps$, $\eta$ and of the amplitude''; both are contradicted by section R1 (amplitude-dependent, $1.3$--$1.75$) and R0 (retracted). The table of ``caveats resolved'' still says the $K=2$ bias excess ``agrees with $K=1$''. These sentences should be struck.
\item \textbf{Noise convention.} The header defines $\sigma^2(\eps)=r/(2\eff)$ and writes the noise term as $\eta\sigma^2$; the derivation section writes the stationary variance as $\eta\sigma^2/2$. With Prop.~\ref{prop:noise} these agree only if $\sigma_\xi^2=r/\eff$ in the derivation and $\sigma^2=a=r/(2\eff)$ in the header. One convention should be used.
\item \textbf{Path-length bound.} $V_T^2/(2\eta T)$ vs the derivable $V_T^2/(\eta T)$ (Prop.~\ref{prop:pathlen}).
\item \textbf{Jump term units.} ``$JD^2/(\eta T)$'' is correct only in event units and only up to the factor $\frac12$ of Prop.~\ref{prop:tanh} ($\rho D^2/(2\eta r)$ in time units).
\item \textbf{Which regime has been proved vs tested.} \emph{(Resolved.)} Assumption~\ref{ass:M} was originally certified only for $\eps\le10^{-4}$ via the Gamma comparator while all experiments used $\eps\in[0.03,0.3]$; Proposition~\ref{prop:Mexact} now certifies it by exact Fourier quadrature for $\eps\le0.3$, $k_0\in[0.6,4]$ (quadrature, not interval arithmetic). Theorem~\ref{thm:Tprime} (TV with explicit constants) remains informative only for $\eps\lesssim10^{-4}$, but the regret theorem no longer depends on it.
\item \textbf{Headline of the information dichotomy.} The dichotomy (bounded versus divergent \emph{expected} information) was read in the notes as ``consistent estimation possible iff $k<1$''. Theorem~\ref{thm:NC} shows this reading is wrong at $\eps=c/T$: no estimator is consistent for any $k$; the dichotomy describes the \emph{typical} accuracy ($\mathcal I_T$ bounded for $k>1$, central quantiles growing like $T^{1-k}$ for $k<1$), not consistency. The wording in the notes and any paper should be changed accordingly.
\item \textbf{Observation model.} The impossibility (``$\mu$ is not consistently estimable when $k_K>1$'') is a statement about the stationary or conditional experiment. For the empty-start record the baseline information diverges at least like $\ln T$ for every $k$ (Prop.~\ref{prop:cold}); the streaming algorithms of Section~\ref{sec:regret} act in the latter setting. The project decision (C4 uses stationary observation) must be stated explicitly in any paper.
\item \textbf{Target form of the regret.} The original target $(d/2)\log T+C\log(1/\eps)+\text{burn-in}$ is not contradicted. Empirically $\E R_n-\log T_n$ is flat in $\eps$ in time units and $\approx-(d/2)\log(1/\eps)$ in events; Theorem~\ref{thm:R} gives an additive constant $O(\eps^{-2}\ln(1/\eps))$, a proof artefact, not a lower bound.
\item \textbf{Status of the earlier K1 summary.} The statement ``stationary-record impossibility, $K=1$: proved'' should read ``proved modulo Theorem~\ref{thm:Tprime}, whose constant chain is hand-derived but evaluated by floating-point quadrature, and informative only for $\eps\lesssim10^{-4}$''.
\end{enumerate}

\subsection{Status matrix}

\begin{center}\small
\begin{tabular}{@{}p{6.1cm}p{1.3cm}p{6.1cm}@{}}\toprule
result & tag & main caveat\\\midrule
Rank-one Hessian; expected curvature; $\lambda_{\min}$ bound & \Pt{} & stationarity assumed\\
Moments, exponential tail, contraction, exact rate $\gamma$ & \Pt{} & constants carry $\beta_{\max}/\beta_{\min}$\\
$K=1$ Fisher formulas; Theorem~\ref{thm:A} (all three regimes) & \Pt{} & $\eps\to0$ at fixed $\mu,\beta$\\
Cold-start lower bound $\ln T$ & \Pt{} & coefficients conjectural\\
Path-term bound for the ridge pair, any $K$, $k_K>1$ & \Pt{} & does not by itself imply inconsistency: needs $\TV(\pi_\theta,\pi_{\theta'})$ ($K=1$: Thm.~\ref{thm:Tprime}); any-$K$ inconsistency is Thm.~\ref{thm:V}\\
Divergence of expected information, $k_K<1$ & \Pt{} & does not imply consistency\\
Matching upper bound $\mathcal I_T^{-1/2}$ (Thm.~\ref{thm:UB}) and non-consistency at $\eps=c/T$ for every $k$ (Thm.~\ref{thm:NC}) & \Pt{}/\Ps{} & $K=1$, $\alpha$ known; NC needs $\TV(\pi_\theta,\pi_{\theta'})$\\
Efficient information for $K\ge2$, $k_K<1$; $k_K=1$ & \Ot{} & numerics only\\
Spectrum: three tiers & \Ps{} & all tiers bounded above and below (Thm.~\ref{thm:tangent}, Thm.~\ref{thm:tiers}); eigenvector convergence rate $O(\eps/g)$; constants loose; $k_K>1$\\
$\lambda_{\min}(I)$ truncation bound & \Ps{} & Lemma~\ref{lem:moments} constants; $20$--$400\times$ loose\\
Tangent minorant: lower bound for every eigenvalue (Thm.~\ref{thm:tangent}) & \Pt{} & exact moment systems solved in 60-digit arithmetic (not interval arithmetic); within $1.3$--$4\times$ of Monte Carlo\\
Regret reduction (Prop.~\ref{prop:reg}), box, projection ($K=1$) & \Pt{} & a-priori box\\
Conditional regret theorem (corrected) & \Ps{} & leverage assumptions\\
Information concentration, exponential tails & \Ps{} & Assumption~\ref{ass:M}; rate constants $10^5$ off\\
E5 exponential tail for $G_n$ & \Pt{} & given leverage conditions\\
Bootstrap with exponential tails & \Ps{} & as above\\
Expected regret $(\tfrac12+o(1))d\ln T+C$, $K=1$ & \Ps{} & box (cost: Prop.~\ref{prop:boxcost}); Assumption M by Prop.~\ref{prop:Mexact}\\
Expected regret, general $K$ (Thm.~\ref{thm:RK}) & \Ps{} & constants uninformative; numerics $K=2$ only\\
Minimax lower bound (Prop.~\ref{prop:LB}) & \Ps{} & hypothesis (U) uniformity; fixed $\eps$\\
Stationary TV bound with explicit constants & \Ps{} & quadrature not interval arithmetic; $\eps\lesssim10^{-4}$\\
Impossibility, stationary record, $K=1$ & \Ps{} & as above\\
Impossibility, stationary record, any $K$ (Thm.~\ref{thm:V}) & \Ps{} & cluster-representation proof not written; $\ln(1/\delta)$ factor\\
Tracking formulas for the linear filter & \Pt{} & --\\
Tracking theorem for forgetting Newton (Thm.~\ref{thm:track}) & \Ps{} & window hypothesis (W); factor $4$ loose; jump class not tight\\
Amplitude effect: $O(A^2)$ closed form (scalar) & \Pt{} & matrix ($\alpha$ unknown) case numerical skeleton only\\
Tracking: window rule, noise coefficient, transfer function & \Nt{} & $K\le3$, $\eps\ge0.2$\\
Two-stage removal of the box (fixed $\eps$) & \Ot{} & --\\
Regret lower bound with explicit $\eps(n)\to0$ & \Ot{} & none found in the literature either\\
Three-tier-aware preconditioning for $K\ge2$ & \Ot{} & --\\\bottomrule
\end{tabular}
\end{center}

\subsection{Position relative to the literature}
The $K=1$ information structure (bounded / logarithmic / divergent) is the Hawkes prelimit of classical CIR theory: the exact Fisher matrix converges to the CIR matrix of Ben Alaya--Kebaier--Tran (Thm.~2.1) as $\eps\to0$ (relative difference $\approx\eps$) and the critical coefficient $2/(\beta(2\mu-\beta))$ equals their critical-case information rate; the CIR literature treats fixed sub/critical/supercritical regimes, and no local-to-critical result was found. Herrera-Mar\'in (two papers) give a two-point KL bound, the Fisher matrix and a Feller-type remark for the near-critical Hawkes likelihood. Discrete analogues in the literature are conditional/ordinary least squares (nearly critical branching with immigration; Anne--Lu--Yu--Zhou 2026 for local-to-unity affine processes), with no likelihood information bound. No Hawkes-specific regret lower bound was found. The novelty of the programme therefore rests on the online regret result (Theorem~\ref{thm:R}), the streaming impossibility statements (Theorems~\ref{thm:imp1}, \ref{thm:V}) and the tracking study, not on the Fisher dichotomy by itself. Every literature claim above was read at abstract or partial full-text level only.

\subsection{Verification record and complexity}\label{sec:verif}
\textbf{Algebraic identities.} Thirty-four identities behind the lemmas (generator computations for $Z_c$, $Z^p$, $e^{\theta Z}$, $S_iS_jS_l$, $e^{-u\lambda}$, $e^{\theta S}$; the Perron objects $\omega,\gamma,\tilde\nu,\kappa$; Lemma~\ref{lem:H}(i)--(iv); Theorem~\ref{thm:A} step 1; the Fisher reduction, efficient information and ridge Rayleigh quotient; the identity behind Lemma~\ref{lem:C1}; the tanh and sinusoid laws) were verified symbolically with exact derivatives of the generator (\texttt{code/verify\_symbolic.py}). \textbf{Inequalities.} Eighteen families (Lemma~\ref{lem:moments} pathwise, Lemma~\ref{lem:expmoment}, Lemma~\ref{lem:G}, the elementary bounds of Lemma~\ref{lem:elem}, Lemma~\ref{lem:gammafact}, Lemma~\ref{lem:F2}, the Bernstein/MGF inequalities, the $\psi$ expansion, the tangent minorant, the Gamma tail, the path-length bounds, the entropy step) had no violation over $10^4$--$10^5$ random cases each (\texttt{code/verify\_inequalities.py}). The probabilistic theorems are chains of these steps with standard tools; they were re-derived step by step but cannot be machine-checked, and Theorem~\ref{thm:V} and the input (H3) of Theorem~\ref{thm:R} are the best targets for independent review.

\textbf{Complexity.} The estimator needs $O(K^2)$ memory and $O(K)$ state; per event the quadrature costs $O(qK^2)$ and the Newton direction $O(K^3)$ (or $O(qK^2)$ with rank-one updates of $F^{-1}$, which pays only for $K\gtrsim100$ in NumPy). The dense $F$-projection, needed by the regret proof, was improved from $\approx55$--$90$ box-QP solves (bisection on the budget multiplier) to $2$--$7$ (safeguarded Newton on the multiplier): $15$--$30\times$ faster, same solution to $1.7\times10^{-13}$. The moment systems of order $2$ and $3$ are Kronecker sums of $A=\mathrm{diag}(\beta)-\mathbf 1\alpha\T$ and are solved by diagonalising $A$: $O(K^4)$ time and $O(K^3)$ memory instead of $O(K^9)$ and $O(K^6)$ ($K=40$ in $21$ ms where the dense system would need $33$~GB). The tangent bound is concave in $t=1/\lambda_0$ ($N(t)=2tM-t^2M_3$ is operator-concave), so a Brent search with $13$--$26$ evaluations replaces a $75$-point grid. See \texttt{verification\_and\_complexity.md}.

\subsection{Open problems, in order of value}
\begin{enumerate}[nosep]
\item Two-stage removal of the box for fixed $\eps$; uniformity (U) for the lower bound; an $\eps$-explicit lower bound for $\eps(n)\to0$.
\item $K\ge2$: usable constants (exploit the three-tier structure; the present $C_K$ is uninformative), $K\ge3$ and $\eps<0.3$ regret numerics.
\item Interval-arithmetic evaluation of $\Phi$, and an analytic Hessian bound for the two-Gamma KL.
\item Theorems~\ref{thm:UB}--\ref{thm:NC} for $K\ge2$ and for $\alpha$ unknown; the regime $c=\eps T\to\infty$ where consistency becomes possible; the law of $\mathcal I_T$ (conjecturally $\mathcal I_T/T^{1-k}$ converges to a boundary-local-time functional).
\item Window hypothesis (W) of the tracking theorem; the matrix amplitude-effect formula; jump-class tracking without the $1/\eta$ loss; $K\ge4$; $\eps<0.4$ at $K\ge2$; the Piggott--Solo full text.
\item Sharper constants for Theorem~\ref{thm:tiers} (the bound on $\E[1/\lambda]$ is $\sim18\times$ loose, which limits the eigenvector bound to $\eps\lesssim10^{-3}$).
\end{enumerate}

\subsection{Note on compilation}
No \LaTeX{} engine was available where this file was produced; it was checked for balanced environments and braces with a script, not compiled. If compilation reports an undefined reference, the label is in the corresponding part file (\texttt{part1\_model\_fisher.tex}, \texttt{part2\_regret.tex}, \texttt{part3\_tracking\_review.tex}).
